\subsection{XGBoost Statistical Baseline}
\label{sec:xgboost_baseline}

A statistical baseline was constructed using the XGBoost gradient-boosted
tree classifier \citep{chen2016xgboost}, which has been applied to lightning
nowcasting in prior work \citep{mostajabi2019nowcasting,holzmann2018error}.
XGBoost was preferred over a convolutional neural network (CNN) because the
available dataset - two seasons of Israeli Lightning Detection Network (ILDN)
observations - is insufficient for reliable deep learning training; CNN-based
nowcasting architectures typically benefit from multi-year dense training
datasets for robust generalisation \citep{leinonen2022seamless}, and
training on two seasons alone would risk severe overfitting.
XGBoost generalises reliably from compact tabular datasets and has been
established as a standard ML approach for lightning nowcasting
\citep{mostajabi2019nowcasting}, providing a useful reference point for
comparison, though differences in study domain, dataset size, and evaluation
protocol limit strict commensurability.

ILDN stroke observations were binned onto a $0.16^{\circ}$ spatial grid
(${\approx}17$\,km) over the study domain
($29$-$35^{\circ}$\,N,\;$31$-$37^{\circ}$\,E) at a temporal resolution
of $\Delta t = 5$\,min, consistent with operational lightning warning systems
\citep{leinonen2022seamless,holzmann2018error} which is expected to capture the dominant convective cell
timescales over the study region.
Both training and evaluation datasets were restricted to November-April,
the months during which convective thunderstorm activity occurs over Israel,
driven by Cyprus Low and Red Sea Trough synoptic systems
\citep{altaratz2003lightning,benami2015lightning}.
The model was trained on November-April 2022-2023 and evaluated on
November-April 2023-2024.

At each forecast origin $t$, features were computed per grid cell spanning
six groups: (i) six multi-scale history sums
$H_w = \sum_{k=0}^{w/\Delta t-1} c(\mathbf{x},t{-}k)$ for
$w \in \{10,20,40,120,240,360\}$\,min \citep{foresti2019using};
(ii) spatial neighbourhood sums over $3{\times}3$ and $5{\times}5$ windows
at both $t$ and $t{-}1$ \citep{holzmann2018error};
(iii) an exponential decay accumulation
$D = \sum_k c(\mathbf{x},t{-}k)\,e^{-0.1k}$ and a persistence ratio
$c / (H_{120}/24 + \varepsilon)$;
(iv) diurnal and seasonal sine-cosine encodings, capturing the morning peak
under Red Sea Trough conditions and the nocturnal maximum of Cyprus Low
frontal systems \citep{benami2015lightning,altaratz2003lightning};
and (v) five static spatial features (normalised latitude, longitude,
land/sea mask, coastline distance, and terrain proxy) encoding
geographically-varying climatology \citep{altaratz2003lightning}.

The binary target for model $\tau$ is defined as
$y_\tau = \mathbf{1}[\sum_{k=\tau/\Delta t - 11}^{\tau/\Delta t}
c(\mathbf{x},t{+}k)>0]$, i.e.\ whether any strike occurs in the specific
one-hour window ending at $t{+}\tau$.
This produces six non-overlapping prediction targets, one per lead time
$\tau \in \{60,120,180,240,300,360\}$\,min, each covering a distinct
future hour.
Six independent XGBoost models were trained, sharing a common feature
matrix $\mathbf{X}$.
Negative training samples were downsampled to 15\% and early stopping was
optimised directly on AUC-PR using a temporally ordered 15\% validation
split to prevent future leakage.
Post-hoc isotonic regression calibration \citep{zadrozny2002transforming}
was applied to correct poorly calibrated scores \citep{caruana2005obtaining}.
Test origins were separated by exactly $\tau$ to ensure independent
verification windows \citep{murphy1993good}.
Figure \ref{fig:xgb_case_study} shows a representative case study from
31 January 2024 illustrating the model's predicted probability field
alongside the observed ground truth at each lead time.

The area under the precision-recall curve (AUC-PR) is used as the primary
metric because AUC-ROC is misleading under severe class imbalance
\citep{davis2006relationship}, supplemented by AUC-ROC and Brier score
\citep{murphy1973hedging} with precision, recall, and F1 at the F1-optimal
threshold.
Precision-Recall curves for each lead time are shown in
Figure \ref{fig:xgb_pr_auc}, and the full suite of skill scores as a
function of lead time is presented in Figure \ref{fig:xgb_scalar_metrics}.

% -
% XGBoost Baseline Figures
% -

\begin{figure*}[t]
    \centering
    \includegraphics[width=\textwidth]{ildn_xgboost_case_study.jpg}
    \caption{XGBoost baseline case study — forecast origin 2024-01-31 09:00 UTC.
    \textbf{Top row:} raw ILDN strike observations in each one-hour
    verification window (strike count shown above each panel).
    \textbf{Middle row:} binary ground truth map on a $0.16^{\circ}$ grid,
    with the number of active cells noted.
    \textbf{Bottom row:} predicted probability field $P(\text{lightning})$
    at lead times 1-6 hours.}
    \label{fig:xgb_case_study}
\end{figure*}

\begin{figure}[H]
    \centering
    \includegraphics[width=\linewidth]{ildn_xgboost_pr_auc.jpg}
    \caption{Precision-Recall curves for the XGBoost baseline at each lead
    time (1-6 hours), evaluated on ILDN data from November-April 2023-2024.
    Diamond markers indicate the F1-optimal operating point.
    AUC-PR values are reported in the legend.}
    \label{fig:xgb_pr_auc}
\end{figure}

\begin{figure}[H]
    \centering
    \includegraphics[width=\linewidth]{ildn_xgboost_scalar_metrics.jpg}
    \caption{XGBoost baseline skill scores as a function of lead time
    (1-6 hours): AUC-ROC, AUC-PR, F1 (at F1-optimal threshold), Brier score,
    Precision, and Recall, evaluated on ILDN data from November-April 2023-2024.}
    \label{fig:xgb_scalar_metrics}
\end{figure}